\begin{exercise}[The two chiral rotation algebras]{I-CH05-001}
Use the mostly-plus metric $\eta=\operatorname{diag}(-,+,+,+)$ and the Lorentz algebra
\begin{equation*}
  [J_{\mu\nu},J_{\rho\sigma}]
  =\ii\left(
    \eta_{\mu\rho}J_{\nu\sigma}
    -\eta_{\nu\rho}J_{\mu\sigma}
    +\eta_{\nu\sigma}J_{\mu\rho}
    -\eta_{\mu\sigma}J_{\nu\rho}
  \right).
\end{equation*}
Put $K_i=J_{0i}$ and
\begin{equation*}
  R_k=\frac12\epsilon_{kij}J_{ij},
  \qquad \epsilon_{123}=1.
\end{equation*}
Interpret Banks's expression $\epsilon_{ijk}J_{ij}$ by counting each
independent antisymmetric spatial pair once, so it denotes $R_k$.  With a sum
over all ordered values of $i,j$, the displayed factor $\frac12$ supplies the
same convention.

Derive the commutators among $R_i$ and $K_i$.  For
$C_i^\pm=R_i\pm\ii K_i$, show that the span of the $C_i^+$ commutes with the
span of the $C_i^-$ and that each span closes on itself.  Give a common
normalization for which each set has the standard $\mathfrak{su}(2)$
commutator.  Finally express the physical rotation generators in terms of the
two commuting sets.
\end{exercise}
